\subsection{Engineering evaluation}\label{sec:engineering}
This section transcribes historical measurements from the repository's
performance/storage reports. No new timing campaign was performed for this
manuscript. Each row applies to its reported revision, fixture and hardware;
newer deep/top-$k$ configurations are not implicitly benchmarked by those rows.
These results evaluate implementation cost and equivalence, not learned fitness.
The source reports and commit identifiers are listed in
Supplementary Section~\ref{app:evidence}.

\paragraph{Runtime and numerical equivalence.}
The second native rollout was measured on an ARM64 Mac using CPython 3.9.6.
Both variants enabled native execution, with the previous implementation at
\code{baa5101} and the update recorded in \code{abd310d}. Each preset used seed
42, 48 resource eaters and 12 prey eaters, one 180-step episode, and medians of
three alternating separate-process measurements per variant. Initialization and
verification hashing were outside the timed loop. Recording cases included
ordinary binary/XZ writing. Figure~\ref{fig:engineering} visualizes the measurements; Supplementary Table~\ref{tab:native} gives exact values.\begin{figure}[p]
\centering
\begin{tikzpicture}
\begin{axis}[
 width=.82\linewidth,height=8.5cm,xmode=log,xmin=.15,xmax=40,
 y dir=reverse,ymin=-.6,ymax=11.6,ytick={0,...,11},
 yticklabels={DQN,PPO,A2C,REINFORCE,GRPO + DQN,Recurrent PPO,QR-DQN,Maskable PPO,MAPPO,Larger-predator PPO,Valley + recording,Flat PPO + recording},
 xtick={.2,.5,1,2,5,10,20},xticklabels={0.2,0.5,1,2,5,10,20},
 xlabel={Loop time (s; logarithmic axis)},tick label style={font=\footnotesize},
 grid=major,grid style={black!10},legend style={font=\footnotesize,at={(.98,.97)},anchor=north east},
 title={\textbf{a}\quad Native-to-native runtime comparison},title style={font=\small},
]
\addplot[only marks,mark=square*,mark size=2.2pt,black!55] table[col sep=comma,x=before,y expr=\coordindex]{figures/native-timings.csv};
\addlegendentry{Previous native}
\addplot[only marks,mark=*,mark size=2.2pt,blue!65!black] table[col sep=comma,x=after,y expr=\coordindex]{figures/native-timings.csv};
\addlegendentry{Updated native}
\end{axis}
\end{tikzpicture}

\vspace{7mm}
\begin{tikzpicture}
\begin{axis}[
 width=.82\linewidth,height=5.3cm,xmin=0,xmax=110,ymin=-.6,ymax=3.6,y dir=reverse,
 ytick={0,1,2,3},yticklabels={401-frame replay,PPO checkpoint,20-run A2C conversion,Deep champion archive},
 xtick={0,25,50,75,100},xlabel={Bytes retained (\% of each fixture's own baseline)},
 tick label style={font=\footnotesize},xmajorgrids,grid style={black!10},
 title={\textbf{b}\quad Independent historical storage fixtures},title style={font=\small},
]
\addplot[xbar,bar width=11pt,fill=black!10,draw=none,bar shift=0pt] coordinates {(100,0)(100,1)(100,2)(100,3)};
\addplot[xbar,bar width=11pt,fill=blue!55,draw=none,bar shift=0pt,nodes near coords,nodes near coords align={horizontal},every node near coord/.append style={font=\footnotesize},point meta=explicit] table[col sep=comma,x expr=100*\thisrow{after}/\thisrow{before},y expr=\coordindex,meta expr=100*\thisrow{after}/\thisrow{before}]{figures/storage-fixtures.csv};
\end{axis}
\end{tikzpicture}
\caption{\textbf{Recorded engineering evidence, without ecological performance claims.}
\textbf{a}, Medians of three alternating separate-process measurements per
variant, 180 steps, seed 42, 48 resource eaters and 12 prey eaters on an ARM64
Mac with CPython 3.9.6. Initialization and validation are excluded; recording
rows include writing. Both versions use native execution. All 72 samples matched
complete-state hashes within their comparison case. Raw replicate distributions
are unavailable here, so uncertainty bars are omitted.
\textbf{b}, Compressed output as a percentage of each independent input fixture;
gray denotes its own 100\% baseline. The first two fixtures are earlier gzip
milestones; later fixtures concern binary conversion and pooled binary/XZ
weights. These percentages cannot be compounded. Exact bytes, times and source
reports appear in Supplementary Tables~\ref{tab:native} and~\ref{tab:storage} and
Supplementary Section~\ref{app:evidence}.}
\label{fig:engineering}
\end{figure}




All 72 simulation samples matched complete-state hashes within their comparison
case, including world, controller parameters, optimizer/buffers, RNGs and
champions. Both recorded cases matched compressed chunk hashes. A separate
1,500-step PPO recording check also matched final state and all 1,501 recorded
frames; its single timing per version is not used as a repeated benchmark.

The earlier native rollout, evaluated against its preceding FP32 Python
implementation, reported $5.30\times$ for PPO and $12.97\times$ for MAPPO, among
other preset-specific gains. The two campaigns have different baselines and
must not be multiplied into a cumulative speedup. Data-path measurements in the
second campaign reported $3.39\times$ for replay expansion and $1.90\times$ for
dashboard summaries, but only $1.03\times$ for raw binary decoding. These are
fixture loop timings with their own repeat counts, not complete-run latencies.

\paragraph{Storage measurements.}
Supplementary Table~\ref{tab:storage} keeps unrelated fixtures separate. MB means $10^6$ bytes;
percentages are calculated from the exact byte counts in the source reports.
The replay delta/gzip and compact-checkpoint rows are earlier milestones, while
current defaults use binary/XZ. They should not be read as one sequential
conversion of a common dataset.



The A2C conversion verified 480 replay chunks and 20 checkpoints before replacing
500 files. Its replay bytes fell from 69,983,482 to 44,842,236; checkpoint bytes
fell from 32,675,595 to 29,434,732. Readable artifacts were retained. This
illustrates that reductions depend on which component dominates an experiment.

The champion fixture contained 22 criterion artifacts but eight unique model
blobs. Its new total includes indexes, metadata, configurations and 521,548 bytes
of weights. A subsequent unchanged archive added 98,993 metadata/index bytes
and zero new weight bytes. These figures concern champion archives, not replay
or full-resume checkpoint storage. Existing results were not silently rewritten.

\paragraph{Memory, correctness and interpretation.}
An earlier 180-step save-path comparison reported peak RSS of 114.53 versus
29.25 MiB and checkpoint writing of 1.04 versus 0.15 seconds. This primarily
reflects streamed checkpoint construction, not an equivalent reduction in live
training memory. Native rollout RSS was not consistently lower: the second
campaign reported approximately 31.9 to 35.2 MiB for recurrent PPO and 34.7 to
33.7 MiB for QR-DQN. Cache/workspace benefits can therefore trade time for memory.

Other recent optimizations were validated by operation counts rather than timed
speed claims: shared configuration serialization once per export identity,
final-winner snapshot copying, deferred closed charts and one score-layout pass
per displayed history sample. Regression checks cover exact continuation,
recording transparency, numerical derivatives, terminal/truncated transitions,
reproductive isolation and Python/JavaScript decoding. A reward-contract review
also checked 216 preset/controller combinations. Passing these checks supports
implementation correctness within their scope; it is not evidence of biological
realism or successful emergent behavior.
